%
% Specifying solver parameters via XML parameter file
%
\chapter{Solver Parameters}
%
%
\section{Interface between CFS++ and OLAS}
%
%
This section deals with specifying solvers and preconditioners for the linear
systems associated with the solution of the different PDEs. It is important to
note at this point that these parameters are of no importance in CFS++ itself,
but will only be passed on to the underlying linear algebra system OLAS. The
necessary interface in CFS++ is proviede by the CFSOLASParams class.

Note that all parameters that can be specified for the individual solvers and
preconditioners are optional. Contrary to the policy in other cases,
default values for missing parameters are currently not specified via the XML
Schema definition and the XMLParamHandler class, but are prescribed in the
SetDefaultParams() methods of the OLAS\_Params class, which constitutes the
method for shuffling parameters from CFS++ to OLAS. See this file or the
doxygen documentation of OLAS for the defaults.
%
%
%
%
\section{Basic specification}
%
%
The PDE section may contain a \xml{solver} element. The element is optional
and can be used to choose a solution algorithm for the linear system
associated to the PDE and in the case of iterative solvers also a
preconditioning method. The two choices are specified as the \xml{type}
and \xml{precond} attributes of the \xml{solver} element. The following two
tables give an overview on the currently available choices for solver and
preconditioner:
%
%
\begin{center}
\begin{tabular}{lp{0.65\textwidth}}
\hline\noalign{\smallskip}
\multicolumn{2}{c}{Currently available choices for type attribute of solver
element}\\
\noalign{\smallskip}\hline\noalign{\smallskip}
Attribute value & solver\\
\noalign{\smallskip}\hline\noalign{\smallskip}
Direct        & direct solution (only supported for LAS)\\
Richardson    & richardson iteration\\
CG	      & conjugate gradient algorithm\\
QMR	      & quasi-minimal residual method\\
hyprePCG      & hypre's version of the classical conjugate gradient algorithm\\
hypreGMRES    & hypre's version of the generalised minimal residual method\\
hypreBICGSTAB & hypre's version of the stabilised bi-conjuigate gradient
                method\\
lapackLU      & direct solution via LAPACK\\
\noalign{\smallskip}\hline
\end{tabular}
\end{center}
%
%
\begin{center}
\begin{tabular}{lp{0.65\textwidth}}
\hline\noalign{\smallskip}
\multicolumn{2}{c}{Currently available choices for precond attribute of solver
element}\\
\noalign{\smallskip}\hline\noalign{\smallskip}
Attribute value & preconditioner\\
\noalign{\smallskip}\hline\noalign{\smallskip}
noPrecond   & no preconditioning (for direct solvers)\\
Id	    & preconditioning with identity matrix, i.e. no preconditioning
              (for iterative solvers)\\
MG	    & proprietory algebraic multigrid method of OLAS\\
Jacobi	    & plain Jacobi preconditioning\\
SSOR	    & symmetric successive over-relaxation\\
ILU	    & incomplete LU decomposition using ILU(0)\\
hypreJacobi & plain Jacobi in hypre\\
hypreAMG    & hypre's algebraic multigrid method\\
hypreILU    & hypre's ILU preconditioner Euclid, ILU(k) currently\\
hypreSPAI   & hypre's SPAI preconditioner ParaSails \\
\noalign{\smallskip}\hline
\end{tabular}
\end{center}
%
%
The default values for \xml{type} and \xml{precond} are \xml{Direct} and
\xml{noPrecond}.
%
%
\section{Parameters for special solvers}
\subsection{CG}
Setting \xml{type} to \xml{CG} selects the native preconditioned conjugate
gradient method implemented of OLAS.
%
%
\begin{center}
\begin{tabular}{lcp{0.5\textwidth}}
\hline\noalign{\smallskip}
element & range of values & meaning \\
\noalign{\smallskip}\hline\noalign{\smallskip}
tol     & $> 0$        & threshold for stopping criterion\\
maxIter & $\mathbb{N}$ & maximal number of iterations\\
\noalign{\smallskip}\hline
\end{tabular}
\end{center}
%
%
\subsection{hypreGMRES}
\subsection{hypreBICGSTAB}
\subsection{hyprePCG}
\subsection{lapackLU}
%
%
\begin{center}
\begin{tabular}{lcp{0.5\textwidth}}
\hline\noalign{\smallskip}
element & range of values & meaning \\
\noalign{\smallskip}\hline\noalign{\smallskip}
tryScaling & yes, no & if set to yes, we let LAPACK try to improve the
                       condition number of the linear system by scaling the
                       rows and columns of the matrix before computing the
                       decomposition\\
refineSol  & yes, no & if set to yes, we let LAPACK perform some steps of
                       iterative refinement after the foward/backward
                       substitution\\
\noalign{\smallskip}\hline
\end{tabular}
\end{center}
%
%
\section{Parameters for special preconditioners}
%
%
\subsection{hypreAMG}
%
%
Setting \xml{precond} to \xml{hypreAMG} selects the BoomerAMG preconditioner
from the hypre library. BoomerAMG is an algebraic multigrid method. The
following table lists the parameters currently supported by the OLAS--hypre
interface. For a detailed description of BoomerAMG see the hypre documentation.
\begin{center}
\begin{tabular}{lcp{0.5\textwidth}}
\hline\noalign{\smallskip}
element & range of values & meaning \\
\noalign{\smallskip}\hline\noalign{\smallskip}
maxLevels & $k\in\mathbb{N}$ & maximal number of levels in multigrid
                               hierarchy\\
alpha     & [0,1]            & central threshold parameter for judging strength
                               of node couplings\\
numSweeps & $k\in\mathbb{N}$ & number of multigrid cycles to perform in each
                               preconditioning step\\
cycleType & V-cycle, W-cycle & specifies the cycling strategy\\
\noalign{\smallskip}\hline
\end{tabular}
\end{center}
%
%
\subsection{hypreILU}
Setting \xml{precond} to \xml{hypreILU} selects the Euclid preconditioner
from the hypre library. Euclid is an incomplete LU decomposition
preconditioner. Currently it only officially supports ILU(k) preconditioning.
\begin{center}
\begin{tabular}{lcp{0.5\textwidth}}
\hline\noalign{\smallskip}
element & range of values & meaning \\
\noalign{\smallskip}\hline\noalign{\smallskip}
level  & $k\in\mathbb{N}$ & generation level of allowed fill-in for ILU(k)
                            approach\\
stats  & yes, no          & print some statistics to standard output\\
memory & yes, no          & print memory some statistics to standard output\\
\noalign{\smallskip}\hline
\end{tabular}
\end{center}
%
%
\subsection{hypreSPAI}
%
%
Setting \xml{precond} to \xml{hypreSPAI} selects the ParaSails preconditioner
from the hypre library. ParaSails can be used as a SPAI or in the symmetric
positive definite case also as FSAI preconditioner. It uses a pre-selected
sparsity pattern for the approximate inverse, which it derives by a powers of
sparsified matrix (PSM) approach.
%
%
\begin{center}
\begin{tabular}{lcp{0.5\textwidth}}
\hline\noalign{\smallskip}
element & range of values & meaning \\
\noalign{\smallskip}\hline\noalign{\smallskip}
thresh & $[-1,1]$ & threshold for determining strong connections in the
                    PSM approach\\
levels & positive integer & number of levels for the PSM approach\\
symmetry & 0, 1, 2 & chooses between SPAI and FSAI\\
loadBalance & yes, no & only important in parallel context\\
reuse & yes, no & do we want to re-use a sparsity pattern\\
\noalign{\smallskip}\hline
\end{tabular}
\end{center}
%
%
\subsection{hypreJacobi}
%
%
This is a simple Jacobi preconditioner. Thus, it does not require or allow the
specification of any parameters.
%
%
\section{Example}
%
%
In this example we tell CFS++ to use the Bi-CGStab solver from the hypre
library together with an SPAI preconditioner:
%
%
\begin{verbatim}
<solver type="hypreBICGSTAB" precond="hypreSPAI">
  <matrix structure="standard" entry="double" storage="sparseNonSym"/>
  <hypreBICGSTAB>
    <tol>     1e-8 </tol>
    <maxIter> 1000 </maxIter>
    <logging>    2 </logging>
    <printLevel> 2 </printLevel>
  </hypreBICGSTAB>
  <hypreSPAI>
    <thresh> -0.4 </thresh>
    <levels>    6 </levels>
    <filter> -0.8 </filter>
    <symmetry>  0 </symmetry>
    <reuse>   yes </reuse>
  </hypreSPAI>
</solver>
\end{verbatim}
